\section{Introduction}

Glycans are essential biological molecules involved in protein folding, cell signaling, immune recognition, and pathogen interaction \citep{varki2017}. Over 50\% of human proteins undergo glycosylation, making accurate glycan structure prediction crucial for understanding biological mechanisms and developing therapeutics \citep{helenius2004}. GlyTouCan, the international glycan structure repository, catalogs over 200,000 unique glycan structures \citep{tiemeyer2016}, highlighting the scale of structural diversity that researchers need to navigate.

AlphaFold 3 (AF3) represents a breakthrough in biomolecular structure prediction, achieving unprecedented accuracy for protein--ligand complexes including glycans \citep{abramson2024}. This advancement has generated considerable excitement in the glycobiology community.

\subsection{The Stereochemistry Problem in AF3 Glycan Modeling}

However, a fundamental challenge emerges when using AF3 for glycan modeling. \citet{huang2025} systematically demonstrated that AF3's standard input formats fail to preserve glycan stereochemistry. Their analysis revealed three input format categories:

\begin{enumerate}
\item \textbf{SMILES format:} AF3 produces incorrect stereoisomers with approximately 62\% stereochemistry accuracy.
\item \textbf{userCCD format:} Partially addresses stereochemistry issues, achieving approximately 82\% accuracy.
\item \textbf{CCD+BAP format:} The only format producing stereochemically correct structures, but requires manual specification taking 30--60 minutes per structure.
\end{enumerate}

This input format bottleneck creates a significant barrier for glycobiology researchers.
